\section{Future Work} \label{chapter:FUTURE_WORK}

Future work is two-fold:
first, improving the performance is desirable to make the solution more practically useful.
Second, additional tool support is required to aid developers in recovery tasks.

\paragraph{Improving the Performance} \label{sec:FUTURE_WORK:1}

Our current implementation introduces a significant execution overhead as presented in Section~\ref{sec:EVALUATION:4}.
The version-aware references resolve to versions of objects dynamically, and even though optimizations such as caching the current versions are possible, most of the current overhead is introduced by the proxies themselves.
Different options could, in parallel, alleviate the performance issues.

Most importantly, the ECMAScript 6 proxy specification is, at the time of this writing, still being draftet.
The proxy implementations in both Chrome and Firefox are only partially integrated with their respective \acp{vm} and \textsc{jit}s.
It seems reasonable to assume that further optimizations will be included in the browsers once the ECMAScript 6 specification is finalized.

Second, we could use proxies less deliberately and exclude some parts of the system from versioning if access to previous states of such parts is not required.
For example, one system part that could be excluded from versioning is the Lively Kernel's \emph{OMeta}~\cite{Warth2007OOL} parser that is used to check for syntax errors.
OMeta creates many objects while parsing code that capture intermediate states of the parser, while in the end only success or failure of the syntax checking needs to be returned.
Given JavaScript's single-threaded, cooperatively scheduled execution it is not possible to switch versions during parsing, so the overhead of the proxies could be removed and direct references be used instead.
A number of modules exist that are often used, but could be considered part of the underlying system and thus not likely to change.

Third, instead of using proxies, version-aware references could be implemented differently.
A similar indirection could be provided using source transformations and ordinary JavaScript functions.
The traps of our proxies could be inserted directly by source transformations around all relevant operations, removing one indirection.

Of these three options, using an alternative to proxies seems most promising.
Using proxies selectively for specific system parts or worlds would only be sufficient if these would not require performance improvements.
Waiting for a faster proxy implementations is an option, but there is at the time of this writing no official release date for the ECMAScript 6 specification.
At the same time, preliminary performance tests indicate that the alternative implementation of version-aware references could be faster.
In particular, microbenchmarks show that going through a function to read a property of an object is only twice as expensive as reading the property directly.

\paragraph{Providing Recovery Tools}

Our implementation allows to preserve and re-establish versions of the Lively Kernel's state.
These versions currently still need to be created explicitly and there are no tools yet to find and manage versions.

Instead of putting the burden of remembering to create a new version on the programmer, the system could create versions of the runtime for any change to an object caused by an action of the programmer, similar to \coexist.
Such actions would entail
\begin{inparaenum}[a)]
    \item manipulating properties of a morph directly with a halo tool or through drag and drop
    \item adding, removing, or editing a script of a morph or a method of a class
    \item evaluating a code snippet
    \item triggering code execution through a mouse or keyboard interaction
\end{inparaenum}
This way, whenever programmers realize changes were inappropriate, they can undo their actions.

The system should support developers in finding and re-establishing relevant states.
Besides preserving versions continuously on a granularity helpful to developers, we want the system to present helpful information to each version.
The system could present three categories of information:
\begin{description}
    \item[when] Versions could be accompanied by a timestamp and be presented in a timelime as in \coexist.
    \item[how] Versions could be annotated with the kind of action that triggered preserving the version such as whether a programmer used a halo button or evaluated a code snippet. This could be supported by recording screenshots or screencasts for versions.
    \item[what] Versions could store information on what was changed between two version: which objects, classes, or files did change, how these objects changed, and how this affected tests and benchmarks
\end{description}

When developers find a relevant previous state, they might want to use it for different purposes:
\begin{description}
    \item[Revisiting previous states] Programmers might want to re-establish a particular state of the system without making changes. For example, they might want to see how an application behaved at a particular moment to compare that to the current state.
    \item[Recovering previous states] Programmers might want to recover state from one version in another version. For example, they could want to recover a particular version of an application or the state of a tool such as a browser.
    \item[Trying alternatives] Programmers might want to try a new idea in an earlier version without loosing neither that version nor any following versions. Therefore, they might want to create a branch as an alternative to the main line of version history.
\end{description}

We want programmers to be able to re-visit versions of the system and to be able to create, merge, and delete lines of history.
Additionally, programmers should be able to copy particular objects from one version to another.

%%% Local Variables:
%%% mode: latex
%%% TeX-master: "../main"
%%% End:
